\subsection{Retrospective finite-path envelope and its interpretation}
\label{sec:v17-predictive-feasibility}
The V16 estimand audit leaves an important empirical question unresolved:
could a correctly implemented, source-mean predictive likelihood
selector match or exceed the source-macro-selected discriminant?
The original Shin2017A result archive alone does not contain the
newly defined selector's source-loss path and cannot establish its
achieved performance. The independent retrospective execution
reported in Section~\ref{sec:v18-fair-predictive} now supplies
such a source-only selection vector. The historical archive also
contains a frozen target AUC path for
every originally \emph{source-feasible} scalar shrinkage candidate
for each outer participant. We use these outcome paths to calculate
an \emph{explicitly post-outcome} feasibility envelope, without
inventing new training results or target-safe parameter choices.

For each target participant $g$, let $\Lambda_g$ denote the original
source-feasible full-domain scalar candidates and $A_g(\lambda)$
the \emph{frozen} held-person AUC corresponding to that participant's
original trained source model. Define the hindsight extremes
\[
 L_g=\min_{\lambda\in\Lambda_g} A_g(\lambda),
 \qquad U_g=\max_{\lambda\in\Lambda_g} A_g(\lambda).
\]
For \emph{any} later selection rule that returns one of the
original $\Lambda_g$ values without altering source training,
representations, model scoring or target trial definitions, its
mean evaluated AUC must satisfy the deterministic inequality
\[
 \frac{1}{G}\sum_gL_g\le
 \frac{1}{G}\sum_gA_g(\widehat\lambda_g)
 \le\frac{1}{G}\sum_gU_g.
\]
This inequality holds regardless of how that rule uses
available \emph{source} data to choose its candidate. It
does \emph{not} cover methods with different representations,
retrained architectures, different target scoring functions,
or a continuous regularisation search outside the archived grid.

For the 29-person Shin2017A study, the original source-macro
selection attained $0.78590$ AUC; the original
validation-class-centred covariance score selector attained
$0.76368$. The deterministic hindsight mean interval
from the 29 frozen paths is \emph{[0.70011, 0.80383]}
AUC. Hence the observed source-macro score is at most
$0.01793$ AUC below a per-participant \emph{target-label-informed
oracle} over the same finite scalar path, while an arbitrary
source-selected future comparator could occupy many different
locations inside that range. Thirteen of the 29 subjects
had an oracle AUC at least 0.01 above the original
source-macro choice.

\begin{table}[tb]
\centering
\small
\caption{Shin2017A original finite-grid AUC feasibility analysis
($G=29$). The min/max envelope uses held-person outcomes after
the experiment; it is not an implementable zero-target predictor
or a confidence interval. No new EEG selector was executed.}
\label{tab:v17-feasible-envelope}
\begin{tabular}{lr}
\toprule
Method or hindsight quantity & Person-macro AUC\\
\midrule
Original source-macro selection & 0.78590\\
Original class-centred covariance selection & 0.76368\\
Minimum per-person frozen-path AUC (hindsight) & 0.70011\\
Maximum per-person frozen-path AUC (hindsight) & 0.80383\\
Oracle mean headroom above source-macro & 0.01793\\
\bottomrule
\end{tabular}
\end{table}

This is an informative \emph{negative identifiability result}:
the archived trials alone cannot establish whether the
new proper predictive-likelihood criterion wins or loses
against source-macro selection. Specifically, source-macro
superiority over the old class-centred covariance score
does not imply superiority over all source-only estimators
within the same original shrinkage family. The newly executed
proper-predictive objective underwent source-only
cross-validation and an immutable 29-choice SHA-256 lock before
joining held-person outcomes, as documented in
Section~\ref{sec:v18-fair-predictive}. The interval in this
section remains a target-label-informed hindsight envelope,
not an achieved model performance interval. Since the source cohort outcomes
have already been inspected, any such run is retrospective,
even with a correctly executed information firewall; a
genuinely independent pre-data cohort would be needed
for new confirmatory inference.
